\documentclass[11pt]{article}
\usepackage[margin=1in]{geometry}
\usepackage{amsmath,amssymb,amsthm,mathtools,bm}
\usepackage{microtype}
\usepackage[T1]{fontenc}
\usepackage{lmodern}
\usepackage{booktabs,array,enumitem}
\usepackage[hidelinks]{hyperref}
\setlist{nosep}
\newtheorem{theorem}{Theorem}[section]
\newtheorem{corollary}[theorem]{Corollary}
\newtheorem{lemma}[theorem]{Lemma}
\newtheorem{proposition}[theorem]{Proposition}
\theoremstyle{definition}
\newtheorem{definition}[theorem]{Definition}
\newtheorem{assumption}[theorem]{Assumption}
\theoremstyle{remark}
\newtheorem{remark}[theorem]{Remark}

\newcommand{\dd}{\mathrm{d}}
\newcommand{\Tr}{\operatorname{Tr}}
\newcommand{\FH}{\mathbf F_{\!H}}

\title{\textbf{Late-Universe Acceleration from Hysterical Response}\\[0.35em]
\large Eventual criticality, vacuum attractor, and a one-parameter background test}
\author{Andrew Bruce Boyd}
\date{Working mathematical formulation --- September 2026}

\begin{document}
\maketitle
\noindent\textit{Computational note.}
Proofs, exploratory experiments, and paper writing were assisted by generative AI.
Mathematical theorems stated in this paper are formally verified in Lean~4 in the
companion specifications, as faithful ports of the TeX arguments at the abstractions
recorded there, except results explicitly tagged as \emph{literature hypotheses}
(standard theorems cited from the literature and not re-proved here) or
\emph{physical inputs} (modeling assumptions outside mathematics).

\begin{abstract}
We develop the late-universe gravitational consequences of the open-system Hysterical Response isolated in the companion response-theory paper.
Homogeneity alone does not force criticality; under a homogeneous FLRW dilution of a conserved density together with a relaxation gap that vanishes as a positive power of that density, the closed-loop gain is nevertheless guaranteed to cross the critical surface at a unique finite scale factor.
Once supercritical, finite-amplitude crossing, a pitchfork attractor, and a minimal fluid completion produce a positive asymptotic density with $w\to-1$ from below.
A slow-pole closure keeps the physical order parameter bounded and yields $\rho_{H,\infty}=m n_c$.
Confronting the adiabatic background with a Planck 2018 flat $\Lambda$CDM baseline shows that late criticality ($z_c\sim1$) produces distance-modulus residuals of several hundredths of a magnitude, while $z_c\gtrsim3$--$5$ approaches $\Lambda$CDM at the millimagnitude level---a one-parameter shape test, not a derivation of the observed magnitude from microscopic coefficients.
\end{abstract}
\tableofcontents
\newpage

\section{Scope and status}
\label{sec:scope}
The response-theory paper~\cite{boyd2026response} isolates the Hysterical Response as the residual component of an already-existing interaction under exponentially stable reduced dynamics.
Galactic modelling continues that residual into an explicit spatial halo~\cite{boyd2026galactic}; the Solar-System null test constrains a local subcritical amplitude~\cite{boyd2026solar}.
The present paper asks a different question: can the same open-system response, once a cosmological channel becomes supercritical under dilution, produce a homogeneous late-time vacuum-like gravitational term?

Four logical layers must be kept separate.
\begin{center}
\begin{tabular}{@{}>{\bfseries}p{0.22\linewidth}p{0.70\linewidth}@{}}
\toprule
Theorem & Consequence of the mathematical assumptions stated here.\\
Modelling assumption & Gap--density class, nonlinear normal form, gravitational closure, or fluid action.\\
Physical input & Expanding FLRW background; non-neutralizing homogeneous mass-energy; conserved number density.\\
Phenomenological confrontation & One-parameter background comparison to a Planck flat $\Lambda$CDM baseline (Sec.~\ref{sec:background}).\\
Not claimed & A microscopic derivation of $K$, $\Gamma_*$, or $C_2$; a full SN/BAO/CMB likelihood; or proof that observed acceleration is Hysterical Response.\\
\bottomrule
\end{tabular}
\end{center}

No new force carrier is introduced.
The scalar order parameters below are coarse-grained coordinates for the gravitational residual response, not fundamental fields.
Homogeneity and isotropy are assumptions of the background model: they license a single background density $n(a)$, not by themselves a critical instability.

\section{Finite-amplitude crossing}
\label{sec:crossing}
Consider a positive scalar response amplitude $X(t)$ obeying the linear unstable equation
\begin{equation}
\label{eq:lin}
\dot X=\lambda(t)X
\end{equation}
with continuous $\lambda$.
Its solution from an initial time $t_i$ is
\begin{equation}
\label{eq:soltime}
X(t)=X_i\exp\!\left(\int_{t_i}^{t}\lambda(s)\,\dd s\right).
\end{equation}

\begin{theorem}[Cosmological finite-amplitude crossing]
\label{thm:crossing}
Let $X:[t_i,t_f]\to(0,\infty)$ solve~\eqref{eq:lin}.
Let $X_*>X_i$.
If
\begin{equation}
\label{eq:threshold}
\int_{t_i}^{t_f}\lambda(s)\,\dd s\ge\ln\frac{X_*}{X_i},
\end{equation}
then there exists $t_*\in[t_i,t_f]$ such that $X(t_*)=X_*$.
If $\lambda(t)>0$ throughout $(t_i,t_f)$ and the inequality in~\eqref{eq:threshold} is strict, the crossing is unique and occurs before $t_f$.
\end{theorem}

\begin{proof}
Equation~\eqref{eq:soltime} gives $X(t_i)=X_i$ and $X(t_f)\ge X_*$.
Continuity of $X$ and the intermediate-value theorem give existence.
If $\lambda>0$, then $\dot X>0$, so $X$ is strictly increasing and the crossing is unique.
Strict inequality places it in the interior.
\end{proof}

\begin{corollary}[Scale-factor form]
\label{cor:scale-factor}
For an expanding background with $H=\dot a/a>0$,
\begin{equation}
\label{eq:sola}
X(a)=X_i\exp\!\left[\int_{a_i}^{a}\frac{\lambda(a')}{a'H(a')}\,\dd a'\right].
\end{equation}
A target $X_*>X_i$ is reached by $a_f$ whenever
\begin{equation}
\label{eq:mastercross}
\boxed{\int_{a_i}^{a_f}\frac{\lambda(a)}{aH(a)}\,\dd a\ge\ln\frac{X_*}{X_i}.}
\end{equation}
\end{corollary}

\begin{remark}[What exponential growth does not prove]
\label{rem:crossing-limits}
The theorem is an existence statement.
It does not explain why a crossing occurs near $a=1$.
That requires independent physics for the onset, rate, seed distribution, and nonlinear saturation.
\end{remark}

For constant positive growth rate $\lambda$, the crossing time is exactly
\begin{equation}
t_*=\frac1\lambda\ln\frac{X_*}{X_i},
\end{equation}
which is the elementary logarithmic threshold formalized in the companion Lean module.

\section{Dilution-driven criticality}
\label{sec:dilution}
The response-theory paper identifies the scalar closed-loop gain $G_H=\chi_R\chi_H$ with criticality at $G_H=1$.
A one-pole slow mode gives schematically $\chi_H\simeq\mathcal R/\gamma_H$.
If the relevant relaxation is collision limited, $\gamma_H\sim n\langle\sigma v\rangle$, and for conserved nonrelativistic number density $n(a)=n_0 a^{-3}$ with slowly varying residue and return susceptibility, one obtains the modelling class
\begin{equation}
\gamma_H\propto a^{-3},\qquad \chi_H\propto a^3,\qquad G_H\propto a^3.
\end{equation}
Thus the cubic law $G=(a/a_c)^3$ is the same closed-loop gain $G_H$ of the response-theory paper, specialized to a density-controlled gap on a homogeneous background---not Newton's constant.

\begin{assumption}[Power-law gap class on a homogeneous background]
\label{ass:eventual}
In an expanding FLRW geometry, a conserved rest-frame density satisfies $n(a)=n_0 a^{-3}$ with $n_0>0$.
On an interval of densities containing the late universe, the slow gap and residue obey
\begin{equation}
\gamma(n)=K n^{\alpha},\qquad K>0,\quad\alpha>0,
\end{equation}
while $\chi_R(n)\mathcal R(n)$ is approximated by a positive constant $\Gamma_*$.
Write $G(n):=\Gamma_*/\bigl(K n^{\alpha}\bigr)$.
\end{assumption}

\begin{theorem}[Eventual dilution criticality]
\label{thm:eventual}
Under Assumption~\ref{ass:eventual}:
\begin{enumerate}
\item $G:(0,\infty)\to(0,\infty)$ is continuous and strictly decreasing, with
$\lim_{n\to0^+}G(n)=\infty$ and $\lim_{n\to\infty}G(n)=0$.
\item There exists a unique critical density $n_c>0$ such that $G(n_c)=1$, namely
\begin{equation}
\boxed{n_c=\bigl(\Gamma_*/K\bigr)^{1/\alpha}.}
\end{equation}
\item Along the expanding FLRW trajectory $n(a)=n_0 a^{-3}$, there is a unique critical scale factor
\begin{equation}
\boxed{a_c=\bigl(n_0/n_c\bigr)^{1/3},}
\end{equation}
with $G\bigl(n(a)\bigr)<1$ for $a<a_c$ and $G\bigl(n(a)\bigr)>1$ for $a>a_c$.
In particular, every expanding homogeneous universe in this class that begins subcritical eventually becomes supercritical at finite $a_c$.
\end{enumerate}
\end{theorem}

\begin{proof}
Continuity and the limits are elementary for $\alpha>0$.
Strict decrease follows from $dG/dn=-\alpha\Gamma_*/(K n^{\alpha+1})<0$.
The intermediate-value theorem gives a unique $n_c$ with $G(n_c)=1$.
Substituting $n=n_0 a^{-3}$ yields $G\propto a^{3\alpha}$, hence a unique $a_c$ as stated.
\end{proof}

\begin{remark}[Homogeneity is not enough]
\label{rem:homo-not-enough}
Theorem~\ref{thm:eventual} uses homogeneity to reduce the background to a single $n(a)$, but the critical crossing is forced by the gap law $\gamma\to0$ as $n\to0$, not by homogeneity alone.
A density-independent gap need never cross $G=1$.
Whether the real gravitational Liouvillian realizes Assumption~\ref{ass:eventual} remains a microscopic input.
\end{remark}

\begin{assumption}[Cubic dilution model]
\label{ass:cubic}
Specialize Theorem~\ref{thm:eventual} to $\alpha=1$, so
\begin{equation}
\label{eq:G}
G(a)=\left(\frac{a}{a_c}\right)^3,
\end{equation}
and, in the supercritical regime $a\ge a_c$, approximate the growing eigenvalue by
\begin{equation}
\label{eq:lambda}
\lambda(a)=\eta H(a)\,[G(a)-1],\qquad \eta>0.
\end{equation}
\end{assumption}

\begin{theorem}[Closed-form dilution-driven growth]
\label{thm:closed}
Under Assumption~\ref{ass:cubic}, a mode with $X(a_c)=X_c>0$ obeys
\begin{equation}
\label{eq:closed}
\boxed{
X(a)=X_c\exp\!\left\{\eta\left[
\frac{a^3-a_c^3}{3a_c^3}-\ln\!\left(\frac{a}{a_c}\right)
\right]\right\},\quad a\ge a_c.}
\end{equation}
In particular, at $a_0=1$,
\begin{equation}
\label{eq:today}
\ln\frac{X_0}{X_c}=\eta\left[
\frac{1-a_c^3}{3a_c^3}-\ln\frac1{a_c}
\right].
\end{equation}
\end{theorem}

\begin{proof}
From $d\ln a/dt=H$ and~\eqref{eq:lambda},
\[
\frac{d\ln X}{d\ln a}=\eta\left[\left(\frac a{a_c}\right)^3-1\right].
\]
Integrating from $a_c$ to $a$ and exponentiating yields~\eqref{eq:closed}.
\end{proof}

\begin{proposition}[Monotonic accumulated growth]
\label{prop:Phi}
Define
\begin{equation}
\Phi(x):=\frac{x^3-1}{3}-\ln x,\qquad x\ge1.
\end{equation}
Then $\Phi(1)=0$ and $\Phi'(x)=x^2-1/x\ge0$, with strict inequality for $x>1$.
Thus the accumulated exponent in~\eqref{eq:closed} grows strictly after criticality.
\end{proposition}

Linear growth cannot continue indefinitely.
A minimal normal form is $\dot X=r(t)X-uX^3$ with $u>0$; for constant $r>0$ the nonzero saturation amplitude is $\sqrt{r/u}$.
The linear crossing theorem is therefore relevant to a target $X_*$ only while $X_*$ remains below saturation.

\section{Nonlinear cosmological attractor}
\label{sec:attractor}
Let $\phi$ be a dimensionless coarse-grained Hysterical order parameter obeying
\begin{equation}
\label{eq:normal}
\dot\phi=r_0[G(a)-1]\phi-u\phi^3,
\qquad r_0>0,\quad u>0,
\end{equation}
with the cubic dilution gain~\eqref{eq:G}.
At fixed $a$, write $r(a)=r_0[G(a)-1]$.

\begin{theorem}[Nonlinear Hysterical branch]
\label{thm:pitchfork}
For fixed $a$, equation~\eqref{eq:normal} has the following equilibria.
\begin{enumerate}
\item If $G(a)\le1$, the only real equilibrium is $\phi=0$; it is linearly stable for $G<1$ and critical at $G=1$.
\item If $G(a)>1$, the neutral equilibrium is unstable and two symmetry-related equilibria exist,
\begin{equation}
\label{eq:branches}
\boxed{\phi_\pm^*(a)=\pm\sqrt{\frac{r_0}{u}[G(a)-1]}.}
\end{equation}
Both nonzero equilibria are linearly asymptotically stable.
\end{enumerate}
\end{theorem}

\begin{proof}
The right-hand side is $f(\phi)=r\phi-u\phi^3=\phi(r-u\phi^2)$.
Thus $\phi=0$ always, and for $r>0$ also $\phi^2=r/u$.
Since $f'(0)=r$, the neutral state changes stability at $r=0$.
At either nonzero equilibrium, $f'(\phi_\pm^*)=r-3u(r/u)=-2r<0$
(literature input: elementary pitchfork linearization; cf.~\cite{perko2001}).
\end{proof}

\begin{remark}[Adiabatic tracking]
\label{rem:adiabatic}
Identifying the cosmological solution with $\phi_\pm^*(a)$ assumes that relaxation toward the stable branch is fast compared with the rate at which $G(a)$ changes.
Near criticality this is least reliable; the full time-dependent equation should be integrated when confronting data.
\end{remark}

\begin{assumption}[Quadratic gravitational response closure]
\label{ass:closure}
On the nonlinear branch, the effective Hysterical gravitational source density is
\begin{equation}
\label{eq:rho-closure}
\boxed{\rho_H=C\rho_m\phi^2,}\qquad C>0,
\end{equation}
where $\rho_m(a)=\rho_{m0}a^{-3}$.
The square makes the macroscopic gravitational source insensitive to the sign of the broken branch.
This is a modelling assumption, not a consequence yet derived from the response-theory Liouvillian.
\end{assumption}

\begin{theorem}[Nonlinear cosmological Hysterical attractor]
\label{thm:attractor}
Assume~\eqref{eq:G}, Theorem~\ref{thm:pitchfork}, Assumption~\ref{ass:closure}, and adiabatic tracking of either stable branch for $a>a_c$.
Then
\begin{equation}
\label{eq:rhoH}
\boxed{
\rho_H^*(a)=C\rho_{m0}\frac{r_0}{u}
\left(a_c^{-3}-a^{-3}\right),\qquad a>a_c.
}
\end{equation}
In particular, $\rho_H^*(a)>0$, is strictly increasing in $a$, and has the finite limit
\begin{equation}
\label{eq:asymptote}
\boxed{\rho_{H,\infty}=C\rho_{m0}\frac{r_0}{u}a_c^{-3}.}
\end{equation}
\end{theorem}

\begin{proof}
From~\eqref{eq:branches}, $\phi_*^2=(r_0/u)[(a/a_c)^3-1]$.
Multiplying by $C\rho_{m0}a^{-3}$ gives~\eqref{eq:rhoH}.
Positivity follows from $a>a_c$; differentiation gives $d\rho_H^*/da=3C\rho_{m0}(r_0/u)a^{-4}>0$.
\end{proof}

\begin{corollary}[Cancellation of dilution]
\label{cor:cancellation}
The large-$a$ constancy is produced by $\rho_m\propto a^{-3}$ and $\phi_*^2\sim G\propto a^3$, so that $\rho_m\phi_*^2$ tends to a constant.
The asymptote is not inserted by hand as a cosmological constant.
\end{corollary}

Define an effective equation-of-state parameter by the bookkeeping identity
\begin{equation}
\label{eq:wdef}
\frac{d\rho_H}{d\ln a}=-3(1+w_H)\rho_H.
\end{equation}

\begin{theorem}[Effective equation of state]
\label{thm:w}
For the adiabatic attractor~\eqref{eq:rhoH},
\begin{equation}
\label{eq:w-attr}
\boxed{
w_H(a)=-1-\frac{1}{G(a)-1}
=-1-\frac{a_c^3}{a^3-a_c^3},\qquad a>a_c.
}
\end{equation}
Consequently $w_H<-1$ at finite $a>a_c$ and $w_H\to-1$ as $a\to\infty$.
\end{theorem}

\begin{proof}
Write $\rho_H=A(a_c^{-3}-a^{-3})$ with $A>0$.
Then $d\rho_H/d\ln a=3Aa^{-3}$, and~\eqref{eq:wdef} yields $1+w_H=-1/(G-1)$.
\end{proof}

\begin{remark}[No fundamental phantom field]
\label{rem:phantom}
The inequality $w_H<-1$ describes an increasing response density approaching a finite asymptote.
It does not by itself imply a fundamental scalar with wrong-sign kinetic term.
\end{remark}

At $a_0=1$, writing $X_0:=\rho_{H0}/\rho_{m0}$ and $B:=C r_0/u$ gives the calibration identity
\begin{equation}
\label{eq:B-cal}
\boxed{B=\frac{X_0}{(1+z_c)^3-1}.}
\end{equation}
A present density ratio therefore fixes only one dimensionless combination once $z_c$ is specified.
For the illustrative fractions $X_0\simeq2.295$, critical redshifts $z_c=2$ and $z_c=3$ give $B\simeq0.088$ and $B\simeq0.036$, with present attractor fractions $26/27$ and $63/64$ and effective $w_{H0}\simeq-1.038$ and $-1.016$ respectively.
These are internal benchmarks, not observational fits.

\section{Critical density and covariant positivity}
\label{sec:covariant}
We now rephrase the same attractor in rest-frame density language and promote it to a minimal covariant fluid.

\begin{assumption}[Power-law dilution class]
\label{ass:powerlaw}
Assumption~\ref{ass:eventual} (already used for eventual criticality), restated for the covariant section.
\end{assumption}

\begin{theorem}[Critical density]
\label{thm:critical-density}
Under Assumption~\ref{ass:powerlaw}, the critical density and redshift are those of Theorem~\ref{thm:eventual}.
For $\alpha=1$,
\begin{equation}
\label{eq:nc-linear}
\boxed{n_c=\frac{\Gamma_*}{K},\qquad G(n)=\frac{n_c}{n}=\left(\frac a{a_c}\right)^3.}
\end{equation}
\end{theorem}

\begin{assumption}[$\mathbb Z_2$ response symmetry]
\label{ass:Z2}
Near the neutral state the deterministic evolution is analytic and invariant under $\phi\mapsto-\phi$.
\end{assumption}

\begin{theorem}[Odd normal form]
\label{thm:odd-normal}
There exist coefficient functions such that
\begin{equation}
\dot\phi=r(n)\phi-u(n)\phi^3+O(\phi^5).
\end{equation}
If the instability is continuous, $r(n_c)=0$.
In the linear-gap class a convenient leading parametrization is
\begin{equation}
\boxed{r(n)=r_0\left(\frac{n_c}{n}-1\right)+\text{higher-order corrections},\qquad r_0>0.}
\end{equation}
If $u_c:=u(n_c)>0$, the supercritical branches satisfy, to leading order,
\begin{equation}
\boxed{\phi_*^2=\frac{r_0}{u_c}\left(\frac{n_c}{n}-1\right).}
\end{equation}
\end{theorem}

\begin{proof}
Oddness removes even powers from the Taylor expansion.
Factoring the cubic truncation gives $\phi[r-u\phi^2]=0$.
For $r>0$ and $u_c>0$, the nonzero roots are real and have linear derivative $-2r<0$.
\end{proof}
\begin{remark}
\label{rem:odd-normal-lean}
Reduced Lean abstraction: the cubic equilibrium algebra is ported; the existence of a $C^\infty$ odd normal form on the response manifold is a modelling/symmetry hypothesis, not a Lean theorem.
\end{remark}

\begin{theorem}[Leading even gravitational response]
\label{thm:even-coupling}
Assume $\rho_H(n,\phi)$ is analytic near $\phi=0$, invariant under $\phi\mapsto-\phi$, and vanishes in the neutral phase.
Then
\begin{equation}
\boxed{\rho_H=C_2(n)\rho_m\phi^2+O(\phi^4),}
\end{equation}
where $C_2$ is dimensionless when $\phi$ is dimensionless.
Symmetry does not determine the sign of $C_2$.
\end{theorem}

\begin{proof}
Evenness removes odd Taylor coefficients; vanishing at $\phi=0$ removes the constant term.
\end{proof}
\begin{remark}
\label{rem:even-coupling-lean}
The leading even monomial is a symmetry argument (TeX); Lean ports the quadratic attractor algebra that follows once this form is assumed.
\end{remark}

Take $\rho_m=mn$ with $m>0$ and approximate $C_2(n)$ by $C_2=C_2(n_c)$ near criticality.

\begin{theorem}[Positive Hysterical attractor density]
\label{thm:positive-attractor}
Under adiabatic tracking of the stable cubic branch,
\begin{equation}
\label{eq:rhoH-n}
\boxed{\rho_H(n)=A(n_c-n),\qquad A:=C_2 m\frac{r_0}{u_c}.}
\end{equation}
Hence
\begin{equation}
\boxed{\rho_{H,\infty}=A n_c=C_2 m\frac{r_0}{u_c}n_c.}
\end{equation}
For linear gap scaling,
\begin{equation}
\boxed{\rho_{H,\infty}=C_2 m\frac{r_0}{u_c}\frac{\Gamma_*}{K}.}
\end{equation}
If $C_2>0$ and $u_c>0$, then $A>0$ and $\rho_H>0$ throughout $0<n<n_c$.
\end{theorem}

\begin{proof}
Insert $\phi_*^2=(r_0/u_c)(n_c/n-1)$ into $C_2 m n\phi_*^2$ and cancel $n$.
\end{proof}

Let $J^\mu=n u^\mu$ be a conserved timelike current with $u^\mu u_\mu=-1$ in $c=1$ units.
Consider the minimal local isotropic adiabatic fluid action
\begin{equation}
\label{eq:action}
\boxed{S_H=-\int d^4x\sqrt{-g}\,\rho_H(n).}
\end{equation}
For such a barotropic current fluid, metric variation yields the standard perfect-fluid identity
\begin{equation}
\label{eq:perfect-fluid}
T_H^{\mu\nu}=(\rho_H+p_H)u^\mu u^\nu+p_H g^{\mu\nu},
\qquad
p_H=n\frac{d\rho_H}{dn}-\rho_H
\end{equation}
(literature hypothesis: variational identity for a barotropic current fluid; see Remark~\ref{rem:fluid}).

\begin{theorem}[Covariant Hysterical stress tensor]
\label{thm:stress}
For $\rho_H=A(n_c-n)$,
\begin{equation}
\boxed{p_H=-A n_c=-\rho_{H,\infty}}
\end{equation}
and
\begin{equation}
\boxed{T_H^{\mu\nu}=-\rho_{H,\infty}g^{\mu\nu}-A n\,u^\mu u^\nu.}
\end{equation}
Thus the homogeneous effective sector is exactly a positive vacuum-like term plus a negative dust-like transient whenever $A>0$.
\end{theorem}

\begin{proof}
$d\rho_H/dn=-A$, hence $p_H=-An-A(n_c-n)=-An_c$.
Substitution into~\eqref{eq:perfect-fluid} gives $\rho_H+p_H=-An$.
\end{proof}

\begin{proposition}[Divergence]
\label{prop:divergence}
For constant $A,n_c$ and conserved $J^\mu$,
\begin{equation}
\boxed{\nabla_\mu T_H^{\mu\nu}=-A n\,a^\nu,\qquad a^\nu:=u^\mu\nabla_\mu u^\nu.}
\end{equation}
Hence the Hysterical sector is separately conserved for geodesic FLRW flow.
\end{proposition}
\begin{remark}
\label{rem:divergence}
Literature/GR identity for the perfect-fluid divergence under $\nabla_\mu(nu^\mu)=0$; not separately Lean-certified here.
\end{remark}

\begin{remark}[Fluid variational identity]
\label{rem:fluid}
Equation~\eqref{eq:perfect-fluid} is part of the effective covariant completion.
Deriving that completion from a Schwinger--Keldysh influence functional remains open.
\end{remark}

\begin{theorem}[Conditional positive cosmological Hysterical attractor]
\label{thm:Lambda}
Assume the hypotheses above, with $m,K,\Gamma_*,r_0>0$, stable saturation $u_c>0$, and positive leading gravitational coefficient $C_2>0$.
For linear density-controlled relaxation,
\begin{equation}
\boxed{\Lambda_H=\frac{8\pi G}{c^2}C_2 m\frac{r_0}{u_c}\frac{\Gamma_*}{K}>0.}
\end{equation}
Moreover, in the broken phase $0<n<n_c$,
\begin{equation}
\boxed{\rho_H+\frac{3p_H}{c^2}=-A(2n_c+n)<0,}
\end{equation}
so its FLRW contribution is accelerating.
\end{theorem}

\begin{corollary}[Effective equation of state in density variables]
\label{cor:w-n}
For $0<n<n_c$,
\begin{equation}
\boxed{w_H(n)=-\frac{n_c}{n_c-n}<-1,}
\end{equation}
with $w_H\to-1$ as $n\to0$.
\end{corollary}

\section{Slow-pole nonlinear closure}
\label{sec:slow-pole}
The Landau branch $\phi_*^2\propto n_c/n-1$ grows without bound as $n\to0$.
A sharper closure places the required low-density enhancement in a slow Liouvillian pole while keeping the physical order parameter bounded.

Let the coarse-grained gravitational response density be
\begin{equation}
\label{eq:response}
\rho_H(n,Y)=-\Tr\!\bigl[\hat\varepsilon_g\,\mathcal L_g(n,Y)^{-1}\mathcal S_g(n,Y)\bigr].
\end{equation}
If $\mathcal L_g$ has a simple slow mode with gap $\gamma(n)$ and pole numerator $\mathcal N(n,Y)$, then
\begin{equation}
\label{eq:pole}
\boxed{\rho_H(n,Y)=\frac{\mathcal N(n,Y)}{\gamma(n)}+\rho_{\mathrm{reg}}(n,Y).}
\end{equation}

\begin{theorem}[Pole-source matching]
\label{thm:pole-match}
Suppose, along a stationary broken branch as $n\to0^+$,
\begin{equation}
\gamma(n)=K_p n^p+o(n^p),\qquad
\mathcal N(n,Y_*)=N_q n^q+o(n^q),
\end{equation}
with $K_p\neq0$, and suppose $\rho_{\mathrm{reg}}\to0$.
Then
\begin{equation}
\rho_H(n,Y_*)=\frac{N_q}{K_p}n^{q-p}[1+o(1)].
\end{equation}
A finite nonzero pole-generated vacuum limit therefore requires $q=p$, in which case
\begin{equation}
\boxed{\rho_{H,\infty}=\frac{N_p}{K_p}.}
\end{equation}
\end{theorem}

\begin{proof}
Divide the two asymptotic expansions in~\eqref{eq:pole}; the regular term vanishes by assumption.
\end{proof}

For linear dilution, $\gamma(n)=Kn+O(n^2)$ and $\mathcal N(n,Y)=N_1 n Y^2+O(n^2)$ with a bounded branch $Y_*^2\to1$ give $\rho_{H,\infty}=N_1/K$.

\begin{assumption}[Two-variable memory closure]
\label{ass:two-var}
Let $Y\in[-1,1]$ be a bounded symmetry-breaking response and $q\in\mathbb R$ an unbounded response-memory amplitude, with
\begin{align}
\label{eq:Ydyn}
\tau\dot Y&=-Y+\tanh q,\\
\label{eq:qdyn}
\dot q&=-Kn q+\Gamma_* Y,
\end{align}
and $\tau,K,\Gamma_*>0$.
\end{assumption}

\begin{theorem}[Neutral-state threshold]
\label{thm:neutral-threshold}
Under Assumption~\ref{ass:two-var}, the neutral state $(Y,q)=(0,0)$ changes linear stability at
\begin{equation}
\label{eq:nc}
\boxed{n_c=\Gamma_*/K.}
\end{equation}
It is linearly stable for $n>n_c$ and unstable for $0<n<n_c$.
\end{theorem}

\begin{proof}
The Jacobian is
\[
J_0=\begin{pmatrix}-1/\tau&1/\tau\\ \Gamma_*&-Kn\end{pmatrix}.
\]
Its trace is negative and its determinant is $(Kn-\Gamma_*)/\tau$.
For a real $2\times2$ system, negative trace and positive determinant imply both eigenvalues have negative real part; a negative determinant gives one positive and one negative real eigenvalue.
\end{proof}

At a stationary point, $q=(n_c/n)Y$, so
\begin{equation}
\label{eq:selfcons}
\boxed{Y=\tanh(GY),\qquad G:=\frac{n_c}{n}.}
\end{equation}
For $G\le1$ the neutral solution is the only stable branch; for $G>1$ the nonzero branches emerge and approach $Y=\pm1$ as $G\to\infty$.
Near criticality,
\begin{equation}
\label{eq:nearY}
\boxed{Y_*^2=\frac{3(G-1)}{G^3}+O((G-1)^2).}
\end{equation}

\begin{assumption}[Reciprocal mass projection]
\label{ass:projection}
In the minimal closure, take
\begin{equation}
\label{eq:projection}
\boxed{\rho_H=m n q Y,}
\end{equation}
where $m>0$ is the conserved gravitating mass per source.
This reciprocity normalization is a physical modelling assumption, not a theorem of the abstract response-theory framework.
\end{assumption}

On a stationary branch,~\eqref{eq:projection} becomes $\rho_H=m n_c Y_*^2$.

\begin{corollary}[Amplitude closure]
\label{cor:amplitude}
Within Assumptions~\ref{ass:two-var}--\ref{ass:projection},
\begin{equation}
\label{eq:rhoInf}
\boxed{\rho_{H,\infty}=m n_c=m\frac{\Gamma_*}{K}.}
\end{equation}
Comparison with the linear-pole vacuum limit gives $N_1=m\Gamma_*$ and residue ratio one.
Near criticality,~\eqref{eq:nearY} implies
\begin{equation}
\rho_H=m n_c\frac{3(G-1)}{G^3}=3m\frac{n_c-n}{G^2},
\end{equation}
so the local slope in $\rho_H=B m(n_c-n)+\cdots$ is $\boxed{B=3}$ as $G\to1^+$ (Lean: algebraic identity plus the $G\to1$ limit).
\end{corollary}

Adopting the constitutive pressure law~\eqref{eq:perfect-fluid}, one finds on a differentiable stable broken branch
\begin{equation}
\label{eq:w-pole}
\boxed{w_H=-1-\frac{2G(1-Y^2)}{1-G(1-Y^2)},}
\end{equation}
with $w_H\to-1$ and $T_H^{\mu\nu}\to-\rho_{H,\infty}c^2 g^{\mu\nu}$ as $n\to0$.
The effective cosmological constant is then
\begin{equation}
\label{eq:Lambda-pole}
\boxed{\Lambda_H=\frac{8\pi G}{c^2}m\frac{\Gamma_*}{K}>0.}
\end{equation}
Within this explicit closure the vacuum density and transition density are no longer independent:
\begin{equation}
\label{eq:master}
\boxed{\rho_{H,\infty}=m n_c=m n_0(1+z_c)^3.}
\end{equation}
Equation~\eqref{eq:master} is a falsifiable relation of the model, not a model-independent prediction of Hysterical Response.
The absolute scale $\Gamma_*/K$ remains an irreducible microscopic input.

\section{Gravity versus electromagnetism}
\label{sec:em}
The generic instability machinery is interaction-agnostic; macroscopic survival is not.

A homogeneous FLRW matter source has positive ordinary mass-energy density.
There is no known ordinary opposite-sign gravitational charge sector that neutralizes it in the manner of electric charge.
Consequently Hysterical neutrality does not force the homogeneous gravitational response to vanish merely because the Universe is homogeneous.

Electromagnetism is different.
Let $Q(t)$ denote the spatially homogeneous net electric charge in a comoving volume.
Exact charge conservation implies $\dot Q=0$, hence
\begin{equation}
Q(t_i)=0\quad\Longrightarrow\quad Q(t)=0
\end{equation}
for all later times under charge-conserving dynamics.

\begin{theorem}[No homogeneous electromagnetic charge generation]
\label{thm:no-charge}
In a closed comoving system with exact total electric-charge conservation, a neutral homogeneous initial state remains neutral.
Therefore a Hysterical electromagnetic instability cannot generate a nonzero $k=0$ net charge from $Q=0$ without additional charge-nonconserving physics.
\end{theorem}

\begin{proof}
Charge conservation integrates to $Q(t)=Q(t_i)$.
\end{proof}

Finite-wavelength charge separation remains allowed.
For a frozen-coefficient electromagnetic characteristic polynomial
\begin{equation}
P(\lambda)=(\lambda+\gamma_H)(\lambda^2+2H\lambda+\omega_p^2)-\beta_{\mathrm{em}}\sigma_H\omega_p^2,
\end{equation}
supercritical feedback $\beta_{\mathrm{em}}\sigma_H>\gamma_H$ with $\omega_p^2>0$ makes $P(0)<0$, hence continuity guarantees a positive real root.
That is a finite-$k$ segregation instability compatible with zero total charge; it is not creation of cosmic net charge.

\section{Background confrontation}
\label{sec:background}
The adiabatic attractor~\eqref{eq:rhoH} supplies a definite one-parameter background once the present density fractions are fixed.
On a flat late-time FLRW geometry, neglecting radiation,
\begin{equation}
\label{eq:E}
\boxed{
E^2(a)=\Omega_{m0}a^{-3}
+\Omega_{H0}\frac{\max(a_c^{-3}-a^{-3},0)}{a_c^{-3}-1},
\qquad a_c=(1+z_c)^{-1}.
}
\end{equation}
Here $\Omega_{H0}$ is matched to the observed dark-energy fraction today; $z_c$ is the single shape parameter.
This is a phenomenological confrontation of the \emph{shape} $1-(a_c/a)^3$, not a microscopic prediction of $\Omega_{H0}$.

\begin{theorem}[Present equation of state]
\label{thm:background-diagnostics}
Under the adiabatic attractor of Theorem~\ref{thm:w}, at $a=1$,
\begin{equation}
\boxed{w_{H0}=-1-\frac{1}{(1+z_c)^3-1}.}
\end{equation}
\end{theorem}

\begin{proof}
Immediate from~\eqref{eq:w-attr} at $a=1$.
\end{proof}

The acceleration onset $z_{\mathrm{acc}}$ is defined by $\rho_m+\rho_H(1+3w_H)=0$ on the attractor branch; its numerical values for a Planck-like baseline appear in Table~\ref{tab:background}.

\paragraph{Planck baseline comparison.}
Adopt the flat Planck 2018 baseline $\Omega_{m0}=0.315$ and $\Omega_{H0}=0.685$ as a literature input~\cite{Planck2018}, and compare~\eqref{eq:E} to flat $\Lambda$CDM with the same fractions.
Table~\ref{tab:background} reports $w_{H0}$, the acceleration onset $z_{\mathrm{acc}}$, luminosity-distance modulus residuals $\Delta\mu(z)=5\log_{10}(d_{L}/d_{L}^{(\Lambda)})$ at $z=0.5$ and $z=1$, and the maximum relative Hubble residual on $z\in[0,3]$.

\begin{table}[h]
\centering
\caption{Adiabatic attractor vs Planck 2018 flat $\Lambda$CDM baseline ($\Omega_{m0}=0.315$). Residuals are \emph{not} a likelihood fit; they quantify shape tension at fixed present fractions.}
\label{tab:background}
\begin{tabular}{cccccc}
\toprule
$z_c$ & $w_{H0}$ & $z_{\mathrm{acc}}$ & $\Delta\mu(0.5)$ & $\Delta\mu(1)$ & $\max|\Delta E/E|$\\
\midrule
$1.0$ & $-1.143$ & $0.93$ & $0.033$ & $0.058$ & $0.113$\\
$1.5$ & $-1.068$ & $0.76$ & $0.016$ & $0.027$ & $0.063$\\
$2.0$ & $-1.038$ & $0.70$ & $0.009$ & $0.015$ & $0.038$\\
$3.0$ & $-1.016$ & $0.66$ & $0.004$ & $0.006$ & $0.017$\\
$5.0$ & $-1.005$ & $0.64$ & $0.001$ & $0.002$ & $0.005$\\
$10$ & $-1.001$ & $0.63$ & $<10^{-3}$ & $<10^{-3}$ & $<10^{-3}$\\
\bottomrule
\end{tabular}
\end{table}

Late criticality ($z_c\sim1$) produces several-hundredths-of-a-magnitude distance residuals and a strongly phantom $w_{H0}$---already in tension with a $\Lambda$-like background.
For $z_c\gtrsim3$--$5$ the background approaches $\Lambda$CDM at the millimagnitude level: the one free shape parameter is then pushed toward early criticality, which is precisely the regime in which Theorem~\ref{thm:eventual} says dilution has long since crossed $G=1$.
A full SN+BAO+CMB likelihood (floating $H_0$, early-universe physics, and growth) is deferred; the present table is the sharpest background statement available without introducing additional free functions.

\paragraph{Growth caveat.}
Because $w_H<-1$ at finite $a$, structure-growth and ISW observables can constrain $z_c$ more tightly than distances alone.
Those probes are not evaluated here; they are listed as the next phenomenological gate.

\paragraph{Series coexistence.}
Galactic modelling and the Solar-System null test constrain spatially organized, often subcritical, response amplitudes~\cite{boyd2026galactic,boyd2026solar}.
The homogeneous cosmological channel of this paper is a different sector: a joint posterior over galactic $V_Q$, Solar $V_\odot$, and cosmological $z_c$ is not attempted.
The logical requirement is coexistence, not identity of amplitudes.

\section{Logical status of the results}
\label{sec:status}
\begin{center}
\small
\begin{tabular}{@{}p{0.46\linewidth}p{0.46\linewidth}@{}}
\toprule
\textbf{Quantity / claim} & \textbf{Status}\\
\midrule
Eventual criticality under gap class~\ref{ass:eventual} & Theorem~\ref{thm:eventual} (TeX general~$\alpha$; Lean: linear-gap $n_c$ / cubic $G(a)$).\\
Finite-amplitude crossing~\eqref{eq:threshold} & Theorem~\ref{thm:crossing} (TeX IVT; Lean: constant-rate threshold).\\
Cubic dilution law $G=(a/a_c)^3$ & Specialization $\alpha=1$ of~\ref{ass:eventual} / Ass.~\ref{ass:cubic}.\\
Closed-form $X(a)$ & Theorem~\ref{thm:closed} (TeX calculus).\\
$\Phi>0$ after criticality & Prop.~\ref{prop:Phi} (Lean).\\
Pitchfork branches~\eqref{eq:branches} & Theorem~\ref{thm:pitchfork} (Lean algebra; Perko context).\\
Odd/even leading forms & Theorems~\ref{thm:odd-normal},~\ref{thm:even-coupling} (symmetry; Remarks~\ref{rem:odd-normal-lean},~\ref{rem:even-coupling-lean}).\\
Quadratic closure $\rho_H=C\rho_m\phi^2$ & Modelling assumption~\ref{ass:closure}.\\
Attractor $\rho_H^*$ identity; $w_H<-1$ & Theorems~\ref{thm:attractor},~\ref{thm:w} (Lean algebra).\\
$w_H$ formula from $d\rho_H/d\ln a$ & TeX proof; Lean ports $w_{\mathrm{eff}}<-1$.\\
Critical density $n_c=\Gamma_*/K$ & Theorems~\ref{thm:eventual},~\ref{thm:critical-density} (Lean).\\
Perfect-fluid identity~\eqref{eq:perfect-fluid} & Literature hypothesis (Remark~\ref{rem:fluid}).\\
Stress / $\Lambda_H>0$ under sign assumptions & Theorems~\ref{thm:stress},~\ref{thm:Lambda} (Lean algebra).\\
Divergence identity & Prop.~\ref{prop:divergence} (Remark~\ref{rem:divergence}).\\
Pole matching $q=p\Rightarrow\rho_{H,\infty}=N_p/K_p$ & Theorem~\ref{thm:pole-match} (asymptotic algebra; Lean: matched-power case).\\
Neutral threshold $n_c=\Gamma_*/K$ & Theorem~\ref{thm:neutral-threshold} (TeX Jacobian; Lean: $n_c$ identity).\\
Reciprocal projection $\rho_H=mnqY$ & Physical modelling assumption~\ref{ass:projection}.\\
$\rho_{H,\infty}=mn_c$; local $B=3$ & Cor.~\ref{cor:amplitude} (Lean algebra + $G\to1$ limit).\\
No homogeneous EM charge from $Q=0$ & Theorem~\ref{thm:no-charge} (Lean).\\
Background shape vs Planck baseline & Phenomenological Table~\ref{tab:background} (literature $\Omega_{m0}$).\\
Observed $\rho_\Lambda$ magnitude from $K,\Gamma_*$ & Microscopic target, not derived.\\
SN/BAO/CMB likelihood; growth/ISW & Deferred.\\
Joint galactic / Solar / cosmological posterior & Deferred coexistence programme.\\
\bottomrule
\end{tabular}
\end{center}

\section{Conclusion}
\label{sec:conclusion}
Homogeneous expansion plus a density-vanishing relaxation gap forces a unique criticality boundary at finite $a_c$; it is the gap law, not homogeneity alone, that does the work.
Once supercritical, accumulated growth forces every finite intermediate amplitude to be crossed; a cubic pitchfork plus quadratic gravitational closure converts that instability into a density approaching a positive constant; a minimal fluid completion supplies constant negative pressure; and a slow-pole closure keeps the physical order parameter bounded while fixing $\rho_{H,\infty}=m n_c$ inside the model.
A one-parameter background confrontation against a Planck flat $\Lambda$CDM baseline then shows that late criticality is background-tense, while $z_c\gtrsim3$--$5$ is presently allowed at the millimagnitude level---strengthening the dilution story without pretending to a parameter-free derivation of $\Omega_{H0}$.
The remaining scientific burden is concentrated: derive $K$ and $\Gamma_*$ from a specified cosmological open-system generator, then run a joint expansion+growth likelihood and a coexistence check against the galactic and Solar constraints of the companion papers.

\bibliographystyle{plain}
\bibliography{paper_refs}
\end{document}
